\chapter{Глоссарий}\label{app:glossary}

\begin{description}[style=nextline, leftmargin=1.2em, itemsep=4pt]
  \item[AprilTag] чёрно-белая метка на коврике; камера дрона узнаёт её и определяет позицию.
  \item[Аварийная посадка] красная кнопка внизу окна и клавиша \klavisha{Esc}: дроны на связи получают команду посадки.
  \item[Взвести] третий шаг цепочки; разрешает старт.
  \item[Борт] дрон шоу с номером N (в коде \kod{show.drone(N)}).
  \item[Борты на связи] настоящие дроны, которые приложение видит по Wi-Fi; их показывает вкладка «Борты».
  \item[Высота] расстояние от пола до дрона в сантиметрах. Лимит по умолчанию 150\cm.
  \item[Дрон шоу] дрон из файла шоу с номером \kod{N}; живому дрону соответствует по выбору или по клетке, на которой он стоит.
  \item[Загрузить] второй шаг цепочки; программа дрона уходит в сам дрон.
  \item[Сцена] левая часть окна: коврик с дронами и их путями, вид сверху (2D) или объёмный (3D).
  \item[Клетка] квадрат коврика 40$\times$40\cm; номер = столбец + 12 $\times$ строка.
  \item[Коврик] поле 12$\times$12 клеток (4,8\,м) с метками.
  \item[Лимит высоты] наибольшая разрешённая высота, от 20 до 200\cm.
  \item[Миссия] программа одного дрона в виде набора шагов, который грузится в дрон.
  \item[Область] часть коврика, которую приложение считает рабочей: весь коврик или четверть.
  \item[Партитура] вкладка справа: временная лента, колонка на каждый борт, время сверху вниз, цвет светодиода линией, кружки на границах шагов.
  \item[Проверки] список замечаний симулятора к шоу.
  \item[Четверть] часть коврика 6$\times$6 клеток; коврик сложен в четыре.
  \item[Шаг] одна команда в программе дрона (взлёт, переход, цвет\ldots).
  \item[Шоу] набор дронов с их программами.
  \item[Симулятор] окно Litebee IDE, которое строит карту и партитуру по файлу шоу.
  \item[Старт] четвёртый шаг цепочки; дроны взлетают после подтверждения.
  \item[Статус «только симулятор»] шаг работает на экране, но на живом дроне не проверен.
  \item[Светодиод] цветная лампа на дроне; управляется командой \kod{rgb}.
  \item[Ожидание] \kod{set_z} на той же высоте: дрон просто висит.
  \item[Полуоборот] половина круга; \kod{circle} считает круги полуоборотами.
\end{description}

\chapter{Что ещё не проверено}\label{app:unverified}

Список для тех, кто готовит первый полёт.

\begin{itemize}
  \item Шаг \kod{code_z} в приложении упаковывается для загрузки так же, как в LiteBee Go, но на живом дроне его ещё никто не запускал. Шаг \kod{circle} уже летал.
  \item Шаг \kod{circle_z} в дрон загрузить нельзя.
  \item У \kod{circle} на дроне шаг длится \kod{count} $\times$ секунды, а симулятор считает секунды за весь шаг; направление \kod{'cw'} на коврике в симуляторе угадано.
  \item Посадка в симуляции занимает 3\,с, а в загруженной программе шаг посадки времени не добавляет.
  \item Свои цвета через \kod{(r, g, b)} и \kod{rgb_value} на живых дронах не проверялись: проверены только восемь именованных цветов.
  \item Модель симулятора приблизительная: дрон стартует под первой клеткой, летит по прямой с постоянной скоростью. Точность живого дрона $10$–$15$\cm.
\end{itemize}
